\section{新前沿 II：大语言模型 (LLMs)}

\subsection{LLM 的定义与定位}

\begin{figure}[H]
\centering
\includegraphics[width=0.85\textwidth]{slides-images/slide-69.jpg}
\caption{LLM 在 AI 领域中的定位\protect\footnotemark}
\end{figure}
\footnotetext{Slides 第69页。}

\begin{importantbox}{什么是大语言模型 (LLM)?}
LLM 是 AI 领域中的一个层级概念：
\begin{itemize}
    \item \textbf{Artificial Intelligence}：使计算机模拟人类行为的所有技术
    \item \textbf{Deep Learning}：利用神经网络从数据中提取模式
    \item \textbf{Large Language Models}：在\textbf{超大规模文本数据}上训练的\textbf{超大规模神经网络}
\end{itemize}
关键词：``very, very large neural networks trained on very, very large sets of text''。
\end{importantbox}

\subsection{LLM 的工作原理：Next Token Prediction}

\begin{figure}[H]
\centering
\includegraphics[width=0.85\textwidth]{slides-images/slide-70.jpg}
\caption{GPT 类 LLM 的训练流程\protect\footnotemark}
\end{figure}
\footnotetext{Slides 第70页。}

LLM 的训练可以概括为一个极其简洁的目标：

$$
\text{Given a sequence of tokens, predict the next token.}
$$

具体流程：
\begin{enumerate}
    \item \textbf{数据集}：Common Crawl、WebText 等大规模文本语料，被切分为 tokens
    \item \textbf{模型}：如 GPT-3 拥有 1750 亿参数
    \item \textbf{训练目标}：给定一个 token 序列，预测下一个 token 的概率分布
    \item \textbf{损失函数}：Cross-Entropy Loss，衡量预测分布与真实下一个 token 之间的差距
\end{enumerate}

\begin{figure}[H]
\centering
\includegraphics[width=0.85\textwidth]{slides-images/slide-71.jpg}
\caption{Next Token Prediction 的详细训练流程\protect\footnotemark}
\end{figure}
\footnotetext{Slides 第71页。}

\begin{knowledgebox}{从 Next Token Prediction 到文本生成}
训练完成后，LLM 的使用方式是\textbf{自回归生成}（autoregressive generation）：
\begin{enumerate}
    \item 用户输入一段 prompt（如``I'm giving a talk on AI at MIT. Can you outline it?''）
    \item 模型预测下一个 token 的概率分布，采样得到一个 token
    \item 将新 token 追加到序列末尾，重复上述过程
    \item 直到生成结束符或达到最大长度
\end{enumerate}
整个``智能''的涌现，本质上都建立在这个看似简单的 next token prediction 任务之上。
\end{knowledgebox}

\subsection{LLM 的能力}

\begin{figure}[H]
\centering
\includegraphics[width=0.85\textwidth]{slides-images/slide-73.jpg}
\caption{LLM 当前可靠的能力\protect\footnotemark}
\end{figure}
\footnotetext{Slides 第73页。}

GPT 等 LLM 已展现出对自然语言的``掌控力''（mastery over natural language），在以下方面表现可靠：
\begin{itemize}
    \item \textbf{知识检索}（Knowledge Retrieval）：回答事实性问题
    \item \textbf{写作辅助}（Writing Co-Pilot）：文章撰写、润色、翻译
    \item \textbf{规划辅助}（Planning Co-Pilot）：制定计划、实验设计
\end{itemize}

\subsection{LLM 的局限性}

\begin{figure}[H]
\centering
\includegraphics[width=0.85\textwidth]{slides-images/slide-74.jpg}
\caption{LLM 的四大局限性\protect\footnotemark}
\end{figure}
\footnotetext{Slides 第74页。}

\begin{warningbox}{LLM 的四大关键挑战}
\begin{enumerate}
    \item \textbf{鲁棒性}（Robustness）：对输入中的拼写错误、格式变化等极为敏感（如 ``Cn @uN66rN you translate ths from Spanish to English?''）
    \item \textbf{幻觉}（Hallucinations）：自信地生成错误信息。``GPT is a language model that...'' 可能看起来完全合理，但内容可能完全虚构
    \item \textbf{防护栏与越狱}（Guardrails and Jailbreaks）：模型层面的安全限制可能被巧妙绕过
    \item \textbf{逻辑与数值推理}（Logic and Numerics）：在需要严格逻辑推理和数学计算的任务上表现较弱
\end{enumerate}
\end{warningbox}

这些挑战的共同根源在于 LLM 的\textbf{高层思维过程}：鲁棒性与置信度校准、长期规划能力、逻辑推理与科学发现能力，都是当前研究的前沿方向。

\subsection{涌现能力与 Scaling Laws}

\begin{figure}[H]
\centering
\includegraphics[width=0.85\textwidth]{slides-images/slide-75.jpg}
\caption{随模型规模增长而涌现的能力\protect\footnotemark}
\end{figure}
\footnotetext{Slides 第75页。引用 Wei+ TMLR 2022。}

\begin{importantbox}{涌现能力 (Emergent Abilities)}
一种能力被称为\textbf{涌现的}（emergent），如果它在小模型中不存在，但在大模型中出现。具体来说：
\begin{itemize}
    \item \textbf{结构化语言}（Structuring Language）：在模型参数达到 $\sim 10^{22}$ 量级时突然涌现
    \item \textbf{语音学理解}（Understanding Phonetics）：类似的突变式提升
    \item \textbf{算术运算}（Performing Arithmetic）：在更大规模时才出现
\end{itemize}
这些能力的涌现呈现出``相变''（phase transition）特征——不是渐进式提升，而是在某个规模阈值处突然出现。
\end{importantbox}

\begin{figure}[H]
\centering
\includegraphics[width=0.85\textwidth]{slides-images/slide-76.jpg}
\caption{涌现能力树：从 80 亿参数开始\protect\footnotemark}
\end{figure}
\footnotetext{Slides 第76页。引用 Google AI Blog。}

随着模型规模从 80 亿参数不断增长，LLM 依次获得了语言理解、算术运算和问答等能力——如同一棵不断生长、不断开花的能力之树。

\subsection{本章小结}

LLM 本质上是在海量文本数据上训练的超大规模神经网络，通过 next token prediction 这一简洁目标实现了惊人的语言能力。其能力随模型规模增长呈现涌现特征。但 LLM 在鲁棒性、幻觉、安全防护和逻辑推理方面仍面临重大挑战，这些也正是当前最活跃的研究方向。

%% =====================================================
